\documentclass[tikz,border=0pt]{standalone}
\usepackage{ctex}
\usetikzlibrary{shadows.blur, positioning, backgrounds, shapes.geometric}

% 定义颜色 - 参考了网站的渐变风格
\definecolor{yan-mint}{RGB}{132, 250, 176}   % #84fab0
\definecolor{yan-sky}{RGB}{143, 211, 244}    % #8fd3f4
\definecolor{yan-bg}{RGB}{30, 30, 30}        % 深色背景模拟暗黑模式

\begin{document}
\begin{tikzpicture}[
    font=\sffamily\bfseries,
    logo node/.style={
        circle,
        minimum size=4cm,
        inner sep=0pt,
        text=white,
        align=center,
        blur shadow={shadow blur steps=10, shadow blur extra rounding=2pt}
    }
]

    % 背景圆 - 使用网站风格的渐变色
    \shade[
        left color=yan-mint,
        right color=yan-sky,
        shading angle=135
    ] (0,0) circle (2.2cm);

    % 内部装饰圆环 (模拟光晕)
    \draw[white, opacity=0.3, line width=1.5mm] (0,0) circle (1.9cm);
    
    % 中心正方形容器 (代码风格)
    \node[
        rounded corners=6pt,
        fill=white!95,
        minimum size=3.1cm,
        blur shadow={shadow blur steps=5}
    ] (box) at (0,0) {};

    % 文字 "Yan"
    \node[scale=5] at (0,0.15) {
        \textcolor{yan-mint}{Y}\textcolor{yan-sky}{an}
    };
    
    % 下方的小字 "Blog" - Favicon模式下隐藏以增大主体
    % \node[scale=0.8, text=gray!80, font=\sffamily] at (0,-0.6) {B L O G};

    % 装饰元素 - 四周的小圆点 (科技感) - Favicon模式下隐藏
    % \foreach \angle in {0, 90, 180, 270} {
    %     \fill[white] (\angle:2.6cm) circle (2pt);
    % }
    
    % 装饰线条 (极简风) - Favicon模式下隐藏
    % \draw[white, line width=0.8mm, line cap=round] (-2.6, 0) arc (180:270:2.6cm);
    % \draw[white, line width=0.8mm, line cap=round] (2.6, 0) arc (0:90:2.6cm);

\end{tikzpicture}
\end{document}
